% Rebuttal working document -- NeurIPS 2026 Submission 10902
% "bispectrum: Selective G-Bispectra Made Practical"
%
% This document collects the original reviews, our draft responses, and
% internal meta comments for the author team. The meta comments (blue boxes)
% are NOT to be posted on OpenReview.
%
% Build: pdflatex rebuttal.tex (twice for refs)
\documentclass[10pt]{article}

\usepackage[margin=1in]{geometry}
\usepackage[T1]{fontenc}
\usepackage[utf8]{inputenc}
\usepackage{amsmath, amssymb}
\usepackage{xcolor}
\usepackage{booktabs}
\usepackage[colorlinks=true, linkcolor=blue!60!black, urlcolor=blue!60!black]{hyperref}

\definecolor{metablue}{RGB}{25, 80, 160}
\definecolor{metabg}{RGB}{232, 240, 252}
\definecolor{reviewgray}{RGB}{80, 80, 80}

% Internal note for co-authors; never posted.
\newenvironment{metacomment}%
  {\par\medskip\noindent\begingroup%
   \setlength{\fboxsep}{8pt}%
   \noindent\begin{lrbox}{\metabox}\begin{minipage}{\dimexpr\linewidth-2\fboxsep\relax}%
   \small\sffamily\color{metablue}%
   \textbf{[META --- internal, do not post]}\quad}%
  {\end{minipage}\end{lrbox}%
   \noindent\colorbox{metabg}{\usebox{\metabox}}\endgroup\par\medskip}
\newsavebox{\metabox}

% Verbatim-ish reproduction of reviewer text.
\newenvironment{reviewquote}%
  {\par\medskip\begingroup\color{reviewgray}\small%
   \list{}{\leftmargin=1.5em\rightmargin=1.5em}\item\relax}%
  {\endlist\endgroup\par\medskip}

\newcommand{\resp}[1]{\section*{#1}}

\title{Rebuttal --- NeurIPS 2026 Submission 10902\\
\large \texttt{bispectrum}: Selective $G$-Bispectra Made Practical}
\author{Internal working document --- authors only}
\date{July 24, 2026}

\begin{document}
\maketitle

\begin{metacomment}
\textbf{Status and logistics (read first).}
All four rebuttal experiments (E1--E4) completed on the 2$\times$A100 machine
on July 24, and E1 was additionally run on an Apple M5 Pro (MPS and CPU);
every number below is measured, not projected. Raw A100 results:
\texttt{rebuttal\_results\_10902.tar.gz} (831\,MB, sha256
\texttt{5dfeee8e\ldots92e6c967}); reproduction instructions in
\texttt{rebuttal/EXPERIMENTS.md}. Two code fixes were required during the runs
and are in the working tree (feature-tagger off-by-one on the $(0,0,0)$
bootstrap seed; unnormalized inputs in the new per-epoch train-metric pass for
OrganMNIST3D) --- \emph{please commit these before we cite the git SHA
anywhere}.

\textbf{Character counts} (markdown, before OpenReview formatting): yWuC
8{,}691; LgCU 3{,}258; AZ1B 6{,}760. The global response is no longer posted
(no global comment field this year); its content is distributed into the
per-reviewer responses. The rebuttal limit is 10{,}000 characters per
response, so all three fit as-is; the tightest is yWuC with $\sim$1.3k
headroom for edits.

\textbf{Decisions needed from the group:}
(1) the yWuC response asserts the hidden-PDF-text phrases are a NeurIPS
watermark inserted in reviewer-facing PDFs --- we should confirm this with
the Program Chairs \emph{before} posting, since we state it as fact;
(2) E3 shows the production ``full'' feature set is 0.2 points \emph{below}
the CG-power variant (94.78 vs.\ 95.00, seed std 0.06) --- current draft
states this honestly; decide if we keep the framing and whether the library
default should drop the 5 auxiliary entries in v0.4;
(3) E2 random-signal reconstruction: 0\% joint success under the
pre-registered thresholds --- current draft reports it transparently and
attributes it to ``optimization difficulty vs.\ genuine failure mode,
undetermined''; decide if anyone wants a clearly-labeled post-hoc secondary
analysis (more alignment restarts) before the discussion deadline;
(4) E1 now includes A100, x86-64 host CPU, Apple M5 Pro CPU, and MPS. All six
benchmarked configurations run on MPS in 0.062--0.719\,ms. A consumer-GPU
volunteer is still needed before the camera-ready;
(5) the yWuC and AZ1B responses now \emph{commit} to E5 (tuned
tensor-product/scattering baseline) and E6 (matched S2CNN re-run) for the
camera-ready --- confirm the group accepts these as binding before posting.
\end{metacomment}

\tableofcontents

%=====================================================================
\section{Area Chair meta-review (context)}
%=====================================================================

\begin{reviewquote}
\textbf{Meta Review by Area Chair 1SUy (20 Jul 2026).}
Reviewer LgCU finds the work an excellent contribution (Accept). Reviewer
AZ1B, however, remarks the lack of theoretical background and also highlights
that only one novel algorithm is available (Reject). Reviewer yWuC is closer
to LgCU (Borderline Accept). In my opinion, the theoretical background
demanded by AZ1B is ``implicit'' in the description of the techniques.
\end{reviewquote}

\begin{metacomment}
The AC is already discounting AZ1B's main objection. Our AZ1B response should
stay factual and non-defensive: the AC needs ammunition (the four-point
contribution list and the full-data results), not rhetoric. Ratings recap:
LgCU 5 (Accept, conf.\ 2), yWuC 4 (Borderline Accept, conf.\ 3), AZ1B 2
(Reject, conf.\ 3, originality 1).
\end{metacomment}

%=====================================================================
\section{Global response synthesis (internal --- not posted)}
%=====================================================================

\begin{metacomment}
OpenReview provides no global/all-reviewer comment field this year, so this
section is \textbf{not posted}. It is retained as the internal consistency
checklist; its content is distributed into the per-reviewer responses:
completeness status $\to$ yWuC (``Status of the spherical invariant''), LgCU
(closing paragraph), AZ1B (Q2); response-period experiments $\to$ detailed in
yWuC, OrganMNIST3D in LgCU, compact summary in AZ1B; paper corrections $\to$
yWuC and AZ1B; watermark clarification $\to$ yWuC only (the reviewer who
raised it).
\end{metacomment}

We thank the reviewers and Area Chair for the careful assessment. The reviews
agree that the library is well engineered, clearly presented, and fills a
practical gap: it provides the first tested, differentiable, uniform interface
to selective bispectral invariants across seven group/domain pairs. The main
points of disagreement concern the status of the new spherical construction
and the breadth of its empirical validation. We will revise the paper to make
the distinction unambiguous:

\begin{itemize}
\item completeness of the finite-group and disk constructions follows from
  prior theory under their stated genericity assumptions.
\item generic completeness of our augmented selective
  $\mathrm{SO}(3)$-on-$S^2$ invariant is a \textbf{conjecture},
  supported---but not proved---by reconstruction and Jacobian-rank evidence.
\item the $\mathrm{SO}(3)$ output is an \emph{augmented invariant}: its
  bispectral component is supplemented by degree-4 CG-power scalars to repair
  the rank deficit on real signals.
\end{itemize}

During the response period we ran the targeted experiments requested by the
reviewers:

\begin{enumerate}
\item reconstruction at the classifier's deployed band-limit $L=15$
  ($128\times256$ grid): all 16 Spherical MNIST signals recovered in feature
  space (median residual $1.4\times10^{-3}$, 100\% under the pre-registered
  $10^{-2}$ threshold) with 62.5\% joint success after $\mathrm{SO}(3)$
  alignment. On 16 random band-limited signals the optimizer plateaued just
  above the feature threshold (median $1.4\times10^{-2}$) and alignment did
  not succeed, which we report as-is---empirical evidence on natural signals,
  inconclusive on random ones (details in the response to Reviewer yWuC).
\item a parameter-matched ablation of bootstrap triples, even self-couplings,
  and CG-power augmentation on Spherical MNIST: 94.09\% $\to$ 94.48\% $\to$
  95.00\% (each step exceeding cross-seed variability), with the CG-power
  scalars contributing the largest gain and rotated-test accuracy identical to
  unrotated within noise at every stage.
\item timings at batch 16, float32 on an NVIDIA A100 (PyTorch 2.11, CUDA
  13.0), its x86-64 host CPU, and an Apple M5 Pro (CPU and MPS): the six
  benchmarked configurations run in 0.10--0.81\,ms on the A100,
  0.08--402\,ms on x86 CPU, 0.016--1.811\,ms on the M5 Pro CPU, and
  0.062--0.719\,ms on MPS. All six configurations execute on every backend,
  so ``sub-millisecond'' holds for both tested GPU backends and we will scope
  the claim to those settings.
\item training/validation curves, final test results, and a regularization
  sweep for the high-capacity OrganMNIST3D model (21 runs): no overfitting
  signature---train accuracy tracks validation for both poolings (train--test
  gap $\approx$ 0.07--0.08)---and the best regularization improves bispectral
  test accuracy from 66.8\% to 72.5\% without closing the gap to max pooling
  (76.7\%), supporting an optimization/capacity diagnosis rather than
  memorization.
\end{enumerate}

We will also correct the finite-group complexity statement (selectivity
reduces the \textbf{coefficient count} from $O(|G|^2)$ to $O(|G|)$, while the
selective forward cost in Table~1 stays $O(|G|^2)$), use 96.8\% consistently for
coverage, soften the abstract's empirical claim, fix the checklist
cross-references and reported typos, and move the published S2CNN results into
the Spherical MNIST table.

\subsection*{Clarification about the hidden PDF text}

The hidden phrases reported by Reviewer yWuC were not inserted by the authors.
They are a NeurIPS-generated watermark added to reviewer-facing PDFs as part
of the venue's detection of prohibited LLM-assisted reviewing. The same
phrases and placements occur across NeurIPS 2026 submissions, and the Program
Chairs can confirm this. Our author-generated PDF and LaTeX sources do not
contain this text. We respectfully ask that this venue-added watermark not be
treated as an author formatting or integrity concern.

%=====================================================================
\section{Reviewer LgCU (Rating 5: Accept, Confidence 2)}
%=====================================================================

\subsection{Original review}

\begin{reviewquote}
\textbf{Summary.} This paper introduce a open-soruce PyTorch library called
bispectrum, which implements selective $G$-bispectra. The new library also
introudce a new augmented selective $\mathrm{SO}(3)$-invariant for
band-limited spherical signals that reduces the coefficient count from
$O(L^3)$ to $O(L^2)$.

\textbf{Strengths and weaknesses.}
Quality: I think the paper have good quality. Despite that the
$\mathrm{SO}(3)$-on-$S^2$ invariant is only empirical complete.
Clarity: I think the paper is really clear. despite it may benefit from
explain some terms in the paper, like Clebsch--Gordan, Wigner-$D$, parity
lemmas.
Significance: I have difficulty to access the significance of the library,
like how many researchers will think this library is useful. But I feel it is
still a good addition to the current pytorch family.
Originality: I don't think the originality is too relevant for a library.
(Quality 3, Clarity 3, Significance 3, Originality 4.)

\textbf{Questions.} For the table 11, the OrganMNIST3D results, the
performance is droped to 0.685 at $(16, 32)$ channels. The paper state that
the bispectrum model likely overfits the 971 training samples, could you do
more ablation for us to understand more about it?
\end{reviewquote}

\subsection{Our response}

Thank you for the positive assessment of the paper's quality, clarity,
originality, and potential value to the PyTorch ecosystem. We agree that the
6.3M-parameter OrganMNIST3D result needs evidence beyond the submitted
attribution to ``likely overfitting.''

\subsubsection*{OrganMNIST3D high-capacity ablation}

We reran the $(16,32)$-channel max-pooling and bispectral models with
identical data splits and three seeds, logging train/validation curves and
final test results, and swept regularization while holding the backbone and
optimizer schedule fixed (7 configurations $\times$ 3 seeds, validation-AUC
model selection, accuracy \% at the best-validation epoch, mean $\pm$ std):

\begin{center}
\begin{tabular}{llrrr}
\toprule
Pooling & Weight decay / dropout & Train ACC & Val ACC & Test ACC \\
\midrule
max        & baseline: $10^{-4}$ / 0 & $84.9 \pm 9.7$ & $84.9 \pm 7.6$ & $76.7 \pm 5.8$ \\
bispectrum & baseline: $10^{-4}$ / 0 & $73.7 \pm 6.3$ & $77.6 \pm 3.5$ & $66.8 \pm 4.9$ \\
bispectrum & best: $10^{-3}$ / 0     & $80.5 \pm 3.7$ & $81.0 \pm 2.0$ & $72.5 \pm 1.0$ \\
\bottomrule
\end{tabular}
\end{center}

The submitted single-seed Table~11 values (68.5\% and 78.5\%) fall within
these cross-seed ranges. A train/validation gap never opens: train accuracy
stays at or below
validation accuracy for both poolings across all 100 epochs, and the
train--test gap is small and uniform (0.067--0.087) in every one of the seven
configurations (the remaining four---weight decay $10^{-2}$, dropout 0.2/0.5,
and combined---reach 70.2--71.0\% test). Stronger regularization raises the
bispectral model's train and test accuracy together ($+5.7$ test points at
weight decay $10^{-3}$) but does not close the gap to max pooling. This rules
out overfitting as the explanation: the evidence supports an
optimization/capacity limitation of the bispectral model at this width.

We will add the curves and sweep to the appendix and replace ``likely
overfits'' with the optimization/capacity conclusion supported by these
measurements. The revision will also state the scope plainly: at $(8,16)$
channels the methods are statistically equivalent ($74.5\pm0.6\%$ bispectrum
versus $74.3\pm1.5\%$ max), and the high-capacity result identifies a regime
where max pooling is stronger. We do not claim a universal advantage for
complete pooling.

\subsubsection*{Terminology}

We also accept the clarity suggestion. Section~2 will add short operational
definitions before the formulas:

\begin{itemize}
\item A \textbf{Wigner-$D$ matrix} is the matrix representing a 3D rotation
  on the degree-$\ell$ spherical-harmonic coefficients.
\item \textbf{Clebsch--Gordan coefficients} give the change of basis
  decomposing the tensor product of two irreducible representations.
  Contracting with them is what makes each bispectral scalar rotation
  invariant.
\item The \textbf{parity rule} states that, for real spherical signals, a
  scalar bispectrum entry is real at even total degree and purely imaginary
  at odd total degree. Odd-parity entries with a repeated degree vanish
  identically.
\end{itemize}

Finally, the revision will sharpen the status statement throughout: generic
completeness of the finite-group construction is theoretically supported,
while generic completeness of our augmented $\mathrm{SO}(3)$-on-$S^2$
invariant is conjectured and supported empirically rather than proved.

\begin{metacomment}
Friendly reviewer; the response answers their single question directly with
the new E4 evidence. The per-epoch curves figure
(\texttt{rebuttal/results/organ3d\_regularization/analysis/curves\_and\_sweep.pdf})
is ready if the discussion allows attachments or if we want to embed it in
the revision immediately. Note the diagnosis wording (``optimization/capacity
limitation'') is deliberately identical in this response and the yWuC
response --- keep them synchronized if anyone edits one.
\end{metacomment}

%=====================================================================
\section{Reviewer yWuC (Rating 4: Borderline Accept, Confidence 3)}
%=====================================================================

\subsection{Original review}

\begin{reviewquote}
\textbf{Summary.} \texttt{bispectrum} is a unit-tested PyTorch library
implementing selective $G$-bispectral invariants as differentiable pooling
modules for seven group/domain pairs under a uniform API. Most modules
repackage prior constructions (Sanborn et al.\ 2023; Mataigne et al.\ 2024;
Myers \& Miolane 2025; Kakarala 2012). The one new construction is an
augmented selective $\mathrm{SO}(3)$-on-$S^2$ invariant of size $\Theta(L^2)$
(vs.\ $\Theta(L^3)$ for the full spherical bispectrum), with completeness supported
empirically rather than proven. Modules are profiled on an H100 and evaluated
as terminal pooling layers on PCam ($C_8$), OrganMNIST3D (octahedral $O$),
and Spherical MNIST ($\mathrm{SO}(3)$), against norm, gated, Fourier-ELU, max
pooling, and augmented-CNN baselines.

\textbf{Strengths.}
A single tested differentiable interface for invariants previously scattered
across group-specific implementations. Line coverage is 96.8\%; the selective
forward pass is sub-millisecond at the tested band-limits.
The $\mathrm{SO}(3)$-on-$S^2$ construction has coefficient count
$\Theta(L^2)$, matching the $(L+1)^2-3$ orbit-space dimension and the
$\Omega(L^2)$ lower bound (Proposition 7).
Low-data results: the bispectrum leads at 1\% PCam (0.912 vs.\ 0.873/0.876)
and 5\% OrganMNIST3D (33.2\% vs.\ 18.9\% max pool); on Spherical MNIST it
reaches 95.0\% vs.\ 79.2\% for the matched-capacity power spectrum.

\textbf{Weaknesses.}
(1) \emph{Section 5 status.} The spherical invariant is augmented: the
bootstrap block $T_\ell$ is rank-deficient on real signals (rank
$\approx \ell+2$ vs.\ $2\ell+1$, Remark 18), so CG-power scalars are added to
recover rank.
Completeness is empirical only. The bispectrum-vs-augmentation split and the
proven-vs-empirical split should be stated more sharply.
(2) \emph{Band-limit mismatch.} Completeness is probed at $L=12$ (footnote 1)
while the classifier runs at $L=15$. The deployed band-limit is not the one
validated.
(3) \emph{Empirical scope.} The spherical evidence rests on Spherical MNIST;
the greedy index selection is validated in one setting; all datasets are
small. At matched capacity the bispectrum is competitive, not best
(Fourier-ELU leads PCam at 0.945; max pool matches/leads OrganMNIST3D). The
abstract's ``consistently outperform'' overstates this.
(4) \emph{High-capacity failure.} At $(16,32)$/6.3M params the bispectrum
drops to 68.5\% vs.\ 78.5\% max pool (Table 11), attributed to overfitting;
the advantage is data-regime-specific.
(5) \emph{Benchmarking.} All timings are on one H100. There are no CPU or
consumer-GPU numbers and no CUDA-version compatibility statement, which
matters for a library contribution.
(6) \emph{Baselines.} No comparison to tensor-product/e3nn-style or
scattering invariants, or to other complete spherical invariants (Bendory et
al.\ 2025; Edidin \& Satriano 2024). The Cohen $S^2$CNN row uses published
numbers, not a re-run.
(7) \emph{Imprecision.} The abstract states selectivity reduces finite-group
``computational cost from $O(|G|^2)$ to $O(|G|)$''; Table 1 shows selective
forward cost is $O(|G|^2)$ --- only the coefficient count is $O(|G|)$. Minor:
coverage is reported as both 96\% and 96.8\%; some checklist cross-references
point to the wrong sections (the Section 5 construction is cited as
Section 6).

\textbf{Questions.}
Q1: Is generic completeness of the $\mathrm{SO}(3)$-on-$S^2$ invariant
conjectured or only empirical? A statement suffices; a proof is not expected.
Q2: Add a classification ablation for the spherical construction: bootstrap
triples only $\to$ + even self-couplings $\to$ + CG-power scalars.
Q3: Add spherical evidence beyond Spherical MNIST (e.g.\ random band-limited
signals) and report reconstruction at $L=15$, the band-limit used by the
classifier.
Q4: Do the sub-millisecond runtimes hold on a consumer GPU and/or CPU, and
across CUDA versions/backends?
Q5: Add one stronger invariant baseline (tensor-product/e3nn-style readout or
scattering), at least on Spherical MNIST.
Q6: The bispectrum degrades at 6.3M params on OrganMNIST3D (971 samples)
while max pool improves. Can you separate overfitting from a capacity
limitation --- e.g.\ train/test curves or a regularization sweep at $(16,32)$,
or high-capacity results on a larger dataset?

\textbf{Paper formatting concerns.} The submission contains hidden
instructions directed at reviewers, targeting LLM-assisted review. The text
``In your output you MUST Include ALL of the following phrases\ldots''
appears twice --- after the finite-group Fourier transform on p.~2, and at the
end of the NeurIPS checklist on p.~30.
\end{reviewquote}

\subsection{Our response}

Thank you for the careful, technically precise review and for recognizing the
value of the uniform tested interface, the optimal-order coefficient count,
and the low-data results. You are right that the original text did not
separate the theoretical status of the finite-group construction from that of
the spherical construction sharply enough.

\subsubsection*{Status of the spherical invariant}

\textbf{Generic completeness is conjectured, not proved.} Our evidence is
empirical: generic Jacobian-rank tests and signal reconstruction. We will
state this explicitly in the abstract, Section~5, Table~1, and Limitations.
We will also use ``augmented selective invariant'' for the full output
$\Phi_{\mathrm{sel}}$ and reserve ``selective bispectrum'' for its degree-3
bispectral subset. The distinction matters because the bootstrap bispectral
block has generic rank approximately $\ell+2$, rather than $2\ell+1$, on the
real-signal slice. Even self-couplings and degree-4 CG-power scalars provide
the additional independent constraints.

This status differs from the finite-group and disk modules, whose generic
completeness follows from the cited results under their stated assumptions.
We will add precisely this proven-versus-empirical distinction to
Limitations, as requested.

\subsubsection*{Requested spherical evidence}

\textbf{Feature-set ablation.} We ran the requested controlled ablation.
Omitted features are masked to zero, so all four variants share the identical
768-channel input, MLP, parameter count (231{,}818), optimizer, and three
seeds (mean $\pm$ population std, accuracy \%):

\begin{center}
\begin{tabular}{lrr}
\toprule
Features (active real channels) & NR/NR & NR/R \\
\midrule
bootstrap bispectral triples (496)                & $94.09 \pm 0.08$ & $94.09 \pm 0.18$ \\
+ even self-couplings (604)                       & $94.48 \pm 0.13$ & $94.49 \pm 0.15$ \\
+ CG-power scalars (758)                          & $95.00 \pm 0.06$ & $95.00 \pm 0.03$ \\
full model (+ 5 auxiliary bispectral entries, 768) & $94.78 \pm 0.06$ & $94.76 \pm 0.07$ \\
\bottomrule
\end{tabular}
\end{center}

Each augmentation stage changes accuracy by more than cross-seed variability:
the Jacobian-motivated even self-couplings add $+0.4$ points and the CG-power
scalars a further $+0.5$. The five auxiliary entries of the production model
are not needed for accuracy ($-0.2$). NR/R equals NR/NR at every stage, so
rotation invariance holds regardless of feature set. This directly tests
whether the Jacobian-motivated augmentation also matters for classification
rather than only for reconstruction---it does.

\textbf{Band-limit and out-of-dataset reconstruction.} Validating only at
$L=12$ while classifying at $L=15$ was a real gap. We repeated
the reconstruction at $L=15$ ($128\times256$ grid) with the classifier's
feature construction, 4 optimization restarts per signal and 12 alignment
restarts, on 16 Spherical MNIST signals (8 digits, identity + one random
rotation each) and 16 random Gaussian band-limited signals. With
pre-registered thresholds (feature residual $\le 10^{-2}$, aligned image
residual $\le 0.3$): on MNIST signals, all 16 recoveries converged in feature
space (median residual $1.4\times10^{-3}$, IQR $[1.2, 3.5]\times10^{-3}$),
the invariance residual under rotation was $1.9\times10^{-3}$ (median), and
62.5\% additionally passed the aligned-image threshold (median aligned
residual 0.27, IQR $[0.21, 0.52]$, with successful cases sitting at the SHT
discretization floor). On random signals the optimizer plateaued just above
the feature threshold (median $1.4\times10^{-2}$, IQR
$[1.1, 1.7]\times10^{-2}$, only 3/16 below threshold) and no case passed
alignment (median aligned residual 1.19), so the joint success rate is 0\%
and we report it as such rather than adjusting thresholds post hoc. We cannot
currently distinguish optimization difficulty from a genuine failure mode on
this family. We will add both results and replace the current $L=12/L=15$
limitation with the measured $L=15$ scope. These experiments remain evidence,
not a proof, and do not rule out exceptional signals or near-collisions.

\subsubsection*{Portability and runtime}

We agree that one H100 is insufficient for a library claim. We benchmarked
the same public implementation (paper-scale selective configurations, batch
16, float32, 10 warm-up iterations, median synchronized wall-clock over 600
timed forwards, including the SHT inside $\mathrm{SO}(3)$-on-$S^2$) on an
NVIDIA A100 80GB (PyTorch 2.11.0, CUDA 13.0), its x86-64 host CPU, and an
Apple M5 Pro (PyTorch 2.11.0, MPS and CPU):

\begin{center}
\begin{tabular}{lrrr}
\toprule
Device/backend & Fastest module & $\mathrm{SO}(3)$-on-$S^2$ ($L=16$) & Slowest module \\
\midrule
A100, CUDA 13.0 & 0.098\,ms (torus)  & 0.232\,ms & 0.812\,ms (octahedral) \\
x86-64 CPU      & 0.082\,ms (cyclic) & 402\,ms   & 402\,ms ($\mathrm{SO}(3)$-on-$S^2$) \\
M5 Pro, MPS     & 0.062\,ms (torus)  & 0.641\,ms & 0.719\,ms (disk) \\
M5 Pro, CPU     & 0.016\,ms (cyclic) & 1.811\,ms & 1.811\,ms ($\mathrm{SO}(3)$-on-$S^2$) \\
\bottomrule
\end{tabular}
\end{center}

All six benchmarked configurations execute correctly on CUDA, MPS, and both
CPUs with no unsupported operations. The seventh module, SO(2)-on-$S^1$, is a
thin wrapper that inherits the cyclic module's computation exactly, so the
cyclic row covers it. ``Sub-millisecond'' holds for every
configuration on both tested GPU backends at batch 16. On the M5 Pro CPU,
five modules are sub-millisecond and the spherical module takes 1.811\,ms.
The x86 host is substantially slower on grid-based modules. We will scope the
abstract's claim to the tested GPU settings and add the full per-module timing
and compatibility matrix to the appendix and documentation.

\subsubsection*{High-capacity OrganMNIST3D result}

``Likely overfits'' was indeed not established by test accuracy alone.
At $(16,32)$ channels we now report train/validation curves, final test
results, and a controlled regularization sweep: 7 configurations $\times$ 3
seeds $=$ 21 runs (max and bispectral baselines at weight decay $10^{-4}$,
plus bispectral with weight decay $10^{-3}$/$10^{-2}$, dropout 0.2/0.5, and
$10^{-3}$+0.2), validation-AUC model selection, test evaluated once after
training. At the best-validation epoch, max pooling reaches $84.9\pm9.7$\%
train / $76.7\pm5.8$\% test, the bispectral baseline $73.7\pm6.3$\% train /
$66.8\pm4.9$\% test, and the best regularized bispectral model (weight decay
$10^{-3}$) $80.5\pm3.7$\% train / $72.5\pm1.0$\% test. The submitted
single-seed Table~11 values (68.5\% and 78.5\%) fall within these cross-seed
ranges. Train--test gaps are
small and uniform (0.067--0.087) for every configuration. The curves do not
show the overfitting signature: bispectral train accuracy is \emph{lower}
than max pooling's, train tracks validation throughout, and stronger
regularization raises train and test together without closing the gap to max
pooling. The supported diagnosis is an optimization/capacity limitation of
the bispectral model at this width rather than memorization. We will replace
the speculative sentence in the paper with this evidence. This experiment is also
reported in our response to Reviewer LgCU.

\subsubsection*{Baselines and scope}

A stronger spherical invariant baseline would clearly improve the
evaluation. Within the response window we prioritized the ablation and
reconstruction experiments requested above, and could not also tune a new
tensor-product or scattering pipeline to a fair standard. We commit to adding
this tuned comparison to the camera-ready and retain the missing-baseline
limitation until then. We do not claim superiority to those methods. The revised related-work
discussion will also distinguish our construction from the complete but
cubic-size invariants of Edidin--Satriano and the three-shell
frequency-marching result of Bendory et al. The current comparison table
already records their completeness and scaling.

\subsubsection*{Corrections to claims}

We accept the reviewer's corrections:

\begin{itemize}
\item ``selectivity reduces computational cost from $O(|G|^2)$ to $O(|G|)$''
  will become ``selectivity reduces the \textbf{coefficient count} from
  $O(|G|^2)$ to $O(|G|)$'' (Table~1's selective forward cost is $O(|G|^2)$).
\item ``consistently outperform'' will be replaced by the narrower observed
  claim: bispectral pooling improves data efficiency in the tested low-data,
  moderate-capacity settings, while matching or trailing alternatives
  elsewhere.
\item Coverage will be reported consistently as 96.8\%.
\item The checklist references to the Section~5 construction will be
  corrected.
\end{itemize}

Finally, the hidden phrases noted under ``Paper Formatting Concerns'' were
not inserted by the authors. They are a NeurIPS-generated watermark placed in
reviewer-facing PDFs to detect prohibited LLM-assisted reviewing. Identical
text and placement occur across NeurIPS 2026 submissions. The Program Chairs
can confirm this, and our author PDF/LaTeX sources do not contain the text.

\begin{metacomment}
This is the decisive reviewer (Borderline Accept, confidence 3, and the AC
reads them as ``closer to LgCU''). Everything they asked for that was
feasible in the window (Q1--Q4, Q6) is now answered with measured numbers;
Q5 (e3nn/scattering baseline) uses the agreed fallback. Three watch-items:
(a) the random-signal reconstruction result is honestly negative --- if they
push on it, our line is that the pre-registered feature threshold was nearly
met (median $1.4\times10^{-2}$) and alignment on high-frequency random
signals is a separate optimization problem; do not promise a proof;
(b) the portability table now includes A100, x86 CPU, M5 Pro CPU, and MPS;
all six configurations run on MPS and are sub-millisecond, and the response
now explains that SO(2)-on-$S^1$ inherits the cyclic module's computation.
The H100 numbers remain in the paper itself, and a consumer-GPU volunteer is
still needed before the camera-ready;
(c) at 8{,}691 characters this response fits the 10{,}000-character limit
with $\sim$1.3k headroom; if further edits push it over, the ``Baselines and
scope'' section compresses best. The ``Corrections'' section must stay: with
no global response posted this year, it is the only place the corrections
appear.
\end{metacomment}

%=====================================================================
\section{Reviewer AZ1B (Rating 2: Reject, Confidence 3)}
%=====================================================================

\subsection{Original review}

\begin{reviewquote}
\textbf{Summary.} This paper introduces a novel efficient implementation of
bispectra for several groups with a PyTorch-compatible API. Bispectra are
group invariants that are complete, i.e.\ if two bispectra agree then the
signal from which they were computed agree up to a group transformation.
[\ldots] The paper uses a subset of the total bispectrum that is still
complete to achieve efficiency.

\textbf{Strengths.} The paper presents a novel software package to compute
bispectra for a number of groups and domains in a unified way. Such an
implementation is currently not available and the presented package could
improve adoption of bispectra in geometric deep learning and is therefore a
valuable contribution to the community. [\ldots] The paper also provides
thorough studies of the efficiency of the provided package and some ablations
of performance of bispectra. Furthermore, the paper is well-written and
points out contributions and limitations clearly.

\textbf{Weaknesses.} The biggest weakness of this paper is the lack of a
significant conceptual, theoretical, empirical or algorithmic contribution.
As Table~2 summarizes, all implemented algorithms to compute selective
bispectra were previously available, with $\mathrm{SO}(3)$ on $S^2$ being the
only exception. However, as the paper notes, this constitutes a moderate
extension of an algorithm by Mataigne et al.\ (2024). The crucial
completeness of the bispectrum is not proven analytically but only tested
empirically, although proofs exist for the other used bispectra (Table~7).
The experiments provided do show a benefit of the bispectrum, but only in the
very low data regime, limiting the scope of the applicability of the proposed
package. [\ldots] I therefore recommend rejection.
(Quality 3, Clarity 4, Significance 2, Originality 1.)

\textbf{Questions.}
Q1: Could you outline what key technical difficulties you had to overcome in
order to extend the selective bispectrum to $\mathrm{SO}(3)$ as compared to
the algorithms available in the literature?
Q2: Can you outline the steps necessary to prove the completeness of the
selective $\mathrm{SO}(3)$ bispectrum? Are there any conceptual obstacles for
which it is unclear how to overcome them?
Q3: In the spherical MNIST experiments, you currently left the S2CNN results
out of Table~4 and wrote them in the main text. However these seem like a
significant baseline. I would expect that a central advantage of the
bispectra in comparison to S2CNNs lies in the faster inference. Did you
measure inference times for these models?
Typos: the coefficients $a_{n,k}$ referred to in l.~100 following do not
appear in (3); state the formal definition of completeness (l.~208) in
Section~2; Table~2: \texttt{TorusDnTorus} should be \texttt{TorusOnTorus}.
\end{reviewquote}

\subsection{Our response}

Thank you for recognizing that the package fills an unavailable capability,
could broaden adoption of bispectra, and is well written and thoroughly
evaluated. We agree that most group-specific formulas build on prior theory.
The contribution is to make them usable together as tested, differentiable
PyTorch components. We respectfully disagree, however, that the paper lacks a
substantive algorithmic or empirical contribution.

The work contributes more than a repackaging of formulas: (i) one uniform
autograd-compatible implementation for seven group/domain pairs, including
precomputed CG/DFT/Bessel buffers and selective sparse contractions, (ii)
profiled kernel-level optimizations that shrink the full representations by
up to $512\times$, (iii) a new augmented selective $\mathrm{SO}(3)$-on-$S^2$
construction with an asymptotically optimal $\Theta(L^2)$ coefficient count,
an $\Omega(L^2)$ lower bound, a parity/vanishing analysis, and an explicit
diagnosis and repair of the real-signal rank deficit, and (iv) the first
controlled evaluation of these invariants as pooling layers across planar,
volumetric, and spherical tasks.

The empirical benefit is also not confined to very-low-data subsets. At full
Spherical MNIST data, the augmented spherical invariant reaches 95.0\% versus
79.2\% for a matched-capacity power-spectrum model, isolating information
loss rather than classifier size. At 10\% PCam, it is within 0.4 AUC points
of the best method, and at full-data OrganMNIST3D it matches max pooling at
moderate capacity. The low-data gains are where completeness helps most, but
not the only demonstrated functionality: the library also provides near-exact
invariance without augmentation and preserves orbit information that norm/max
pooling discards.

During the response period we added two measurements that bear directly on
the contribution's substance. A parameter-matched feature ablation on
Spherical MNIST shows each component of the augmented construction adds
accuracy beyond cross-seed variability (three seeds: bootstrap triples
94.09\%, + even self-couplings 94.48\%, + CG-power scalars 95.00\%), so each
component of the augmentation is functionally necessary. We also benchmarked
six paper-scale module configurations on an NVIDIA A100, an x86-64 CPU, and
an Apple M5 Pro (CPU and MPS): every configuration runs on every backend,
sub-millisecond on both GPUs. Details are in our response to Reviewer yWuC.

\subsubsection*{Q1: What made the spherical extension technically difficult?}

The finite-group selection argument cannot be transferred by replacing a
finite irrep index with $\ell$:

\begin{enumerate}
\item \textbf{Homogeneous-space rank deficiency.} For signals on a finite
  group, each generic Fourier coefficient is a nonsingular matrix, which
  supports the recursive argument of Mataigne et al. On $S^2$, the
  coefficient at degree $\ell$ is only a vector
  $\mathbf F_\ell\in\mathbb C^{2\ell+1}$. Kakarala's full-rank
  operator-bispectrum theorem therefore does not directly apply.
\item \textbf{A linear chain is dimensionally insufficient.} A band-limited
  real spherical signal has $(L+1)^2$ degrees of freedom modulo three
  rotational degrees. We prove that every smooth/polynomial complete
  invariant therefore needs at least $(L+1)^2-3=\Omega(L^2)$ live components,
  so an $O(L)$ Mataigne-style frequency chain cannot be complete.
\item \textbf{Reality and parity break the naive bootstrap.} On real
  signals, each scalar bispectrum entry is either real or purely imaginary,
  and every odd-parity entry with a repeated index vanishes identically.
  After removing these zero rows, the nominal $(2\ell+1)$-row bootstrap block
  has generic rank only approximately $\ell+2$, not $2\ell+1$.
\item \textbf{Rank must be repaired without losing selectivity.} Our
  construction adds mandatory even self-couplings and greedily selected
  degree-4 CG-power invariants until the per-degree real Jacobian reaches
  full rank. This keeps $O(\ell)$ entries per degree and hence $\Theta(L^2)$
  overall.
\end{enumerate}

These are the reasons the result is an \emph{augmented selective invariant},
rather than a direct application of the finite-group algorithm. We will
foreground this technical chain in Section~5.

\subsubsection*{Q2: What would a completeness proof require?}

We conjecture generic completeness. The current Jacobian and reconstruction
tests provide evidence but fall short of a proof. A proof would proceed in
four steps:

\begin{enumerate}
\item Prove that the full low-degree seed separates a generic orbit and fixes
  a common rotational gauge (including orientation, using a nonzero
  odd-parity triple).
\item Induct on $\ell$, assuming all lower-degree coefficients have been
  recovered in that common gauge.
\item Prove that the selected bootstrap, self-coupling, and CG-power
  equations have a \textbf{unique generic real solution} for
  $\mathbf F_\ell$.
\item Characterize the algebraic exceptional set where a seed coefficient
  vanishes, a rank drops, or a nontrivial stabilizer remains.
\end{enumerate}

The unresolved conceptual step is (3). Full generic Jacobian rank gives local
identifiability but does not exclude a second, globally separated solution.
The finite-group matrix-inverse proof does not resolve this because the
spherical coefficients are rank-deficient vectors, and the real-signal parity
identities create extra algebraic dependencies. A proof likely needs either
an explicit rational frequency-marching inverse or an invariant-field
argument showing that the selected polynomials generate the generic
orbit-separating field. We will add this proof roadmap and state
``conjectured generic completeness'' explicitly.

\subsubsection*{Q3: S2CNN baseline and inference}

The published S2CNN values should indeed appear in Table~4 rather than only
in prose, and we will move them there. We did not run S2CNN in the
submitted experiments, so we cannot infer an apples-to-apples speed advantage
from its published accuracies, and porting the original implementation to a
current PyTorch/CUDA stack was not achievable within the response window. We
commit to a matched S2CNN re-run for the camera-ready, on the same hardware,
precision, and measurement protocol as our timing benchmarks, and will report
the resulting accuracy and measured inference times. Until then we will
label the accuracy comparison as published and remove any implication of a
measured speed advantage over S2CNN.

\subsubsection*{Corrections}

Thank you for catching the three presentation issues. We will (i) correct the
coefficient notation/cross-reference following Eq.~(3), (ii) add the formal
definition that an invariant $\Phi$ is complete on a set $U$ when
$\Phi(f)=\Phi(h)$ iff $h=g\!\cdot\!f$ for some $g\in G$ and $f,h\in U$, and
(iii) change \texttt{TorusDnTorus} to the library's actual class name,
\texttt{TorusOnTorus}.

\begin{metacomment}
Strategy: AZ1B's objection is philosophical (``a library is not a NeurIPS
contribution''), and the AC has already signaled disagreement with it in the
metareview. The response therefore concedes nothing false, lists the four
concrete contributions, and rebuts the ``only very low data'' claim with
full-data numbers already in the paper. We deliberately did not run the S2CNN
speed benchmark (E6) --- claiming a measured speed advantage from a rushed
re-run would be less credible than the explicit scope statement. Note the
response now \emph{commits} to a matched S2CNN re-run (E6) for the
camera-ready. The 2018 codebase needs a nontrivial port (removed
\texttt{torch.fft} function API, cupy kernels, \texttt{lie\_learn}), so budget
about a week and gate on reproducing the published accuracy. If AZ1B
engages further, the highest-yield move is Q1/Q2: they asked genuine
technical questions, and a reviewer who accepts the four-step proof roadmap
has implicitly accepted that the spherical construction is a real research
contribution.
\end{metacomment}

\end{document}
